
\section{Preliminaries}\label{sec:prelim}
In this section we fix the syntax and semantics of Timed Computation Tree Logic (TCTL) over the non-punctual fragment, the timed transition systems that interpret it, and the zones used as symbolic representations of clock valuations in the tableau construction.



\subsection{Syntax and semantics of TCTL}

TCTL extends branching-time temporal logic—specifically Computation Tree Logic (CTL)—with quantitative time constraints on temporal modalities. 


To reason about time constraints, {clocks} are introduced, which are real-valued variables that increase uniformly over time. A {clock valuation} is a function $\nu : \mathcal{C} \to \mathbb{R}_{\ge 0}$ that assigns a non-negative real value to each clock in a finite set $\mathcal{C}$. We write $\clockValuations{\mathcal{C}}$ for the set of all clock valuations.





Clocks constraints are expressed using time bounds. Time bounds are defined as follows.
\begin{definition}[Time bounds]\label{def:time_bounds}
    A time bound $\kI$ is a non-empty interval of $\mathbb{R}_{\ge 0}$ with {integer} endpoints, i.e., $\kI$ is of the form $[a,b]$, $(a,b]$, $[a,b)$, or $(a,b)$ where $0 \le a < b$, $a \in \Int_{\ge 0}$, and $b \in \Int_{\ge 0} \cup \{\infty\}$.
\end{definition}


We write $<b$ and $\leq b$ as a shorthand for $[0,b)$ and $[0,b]$, as well as $>a$ $\geq a$ as a shorthand for $(a,\infty)$ and $[a,\infty)$, respectively. Note that we allow only non-singular (non-punctual) intervals as time bounds, which means that punctual constraints of the form $x=a$ (equivalently, singular intervals $[a,a]$) are not expressible in our syntax. This restriction is necessary to ensure decidability of the satisfiability problem for TCTL: with punctual constraints satisfiability is undecidable~\cite{AlurCourcoubetisDill1990}, whereas the non-punctual fragment is decidable~\cite{fsttcs2000-LN,LaTorreNapoli2001TimedTreeAutomata}. We emphasise that ``non-punctual'' forbids only singular intervals; {strict} bounds such as $(a,b)$, $(a,b]$, $[a,b)$, $<b$, and $>a$ are allowed.


We build TCTL formulas from a set of atomic propositions $\mathcal{P}$. Atomic propositions are interpreted as Boolean variables that can be true or false in a given state. We can define the syntax of TCTL formulas inductively as follows. 

%Now we can define the syntax of TCTL formulas as follows. 
\begin{definition}[Syntax of TCTL]\label{def:tctl_syntax}
    Given a set of atomic propositions $\mathcal{P}$, a finite set of clocks $\clockSet$ the syntax of TCTL formulas we consider in this work is defined as follows.

\[
\begin{aligned}
\varphi \;::=\;
&\;\top
\mid \bot
\mid p
\mid \neg p
\mid\;
\varphi \land \varphi
\mid \varphi \lor \varphi
\\[2pt]
&\mid\;
\EU{\varphi}{\varphi}
\mid \AU{\varphi}{\varphi}
\mid\;
\tEU{\varphi}{\varphi}{\kI}
\mid \tAU{\varphi}{\varphi}{\kI}
\\[2pt]
&\mid\;
\AR{\varphi}{\varphi}
\mid \ER{\varphi}{\varphi}
\mid\;
\tAR{\varphi}{\varphi}{\kI}
\mid \tER{\varphi}{\varphi}{\kI}.
\end{aligned}
\]

\end{definition}





The operators $\mathbf{U}$ and $\mathbf{R}$ denote the (untimed) {until} and {release} modalities, respectively, while $\mathbf{U}_{\kI}$ and $\mathbf{R}_{\kI}$ denote their time-bounded variants. The untimed variants are the special case $\kI=[0,\infty)$ and are kept as primitives only for readability.

Unlike (untimed) CTL, the logic deliberately includes {no} next-state operator. In the dense-time semantics time elapses continuously, so no configuration has an immediate temporal successor and a ``next'' modality is not definable over the reals~\cite{AlurCourcoubetisDill1990}. The symbols $\EX{}$ and $\AX{}$ that occur in the tableau calculus of Section~\ref{sec:tableau} are therefore {not} object-language operators: they are proof-theoretic markers that denote, respectively, the existential and universal quantification over the {discrete} successors generated by the construction. Their meaning is fixed operationally by the successor-generation rule~\eqref{tag:ex} and the discrete relation $\rightarrow_d$, never by the satisfaction relation $\models$ below. This is the precise sense in which our object logic is the next-free dense-time fragment of~\cite{AlurCourcoubetisDill1990,fsttcs2000-LN}, while CTL-style next reasoning survives only internally to the proof search.
%
%Specifically, timed bounds $\kI$ range over non-empty intervals of $\mathbb{R}_{\ge 0}$ 
%with real endpoints, that is,
%\[
%\kI \in \{\, [a,b], (a,b], [a,b), (a,b) \mid 
%0 \le a < b,\; a \in \Int_{\ge 0}, b \in \Int_{\ge 0} \cup \{\infty\} \,\}.
%\]







We assume formulas are given in {negation normal form} (NNF), where negation may only occur directly in front of atomic propositions. Thus, a subformula $\neg \varphi$ is permitted only if $\varphi$ is an atomic proposition $p \in \mathcal{P}$ (negations of $\top,\bot$ reduce to $\bot,\top$).

The Boolean connectives $\rightarrow$ and $\leftrightarrow$ are not taken as primitive, since they can be defined by:
$ \varphi \rightarrow \psi \;\equiv\; \neg \varphi \lor \psi$ and $\varphi \leftrightarrow \psi \;\equiv\;
(\varphi \land \psi) \lor (\neg \varphi \land \neg \psi).$ Similarly, the standard derived modalities are not included as primitives since they are definable using until and release, i.e. $\mathrm{EF}\,\varphi \;\equiv\; \EU{\top}{\varphi}$, $\mathrm{AF}\,\varphi \;\equiv\; \AU{\top}{\varphi}, $ and, dually,
$\mathrm{EG}\,\varphi \;\equiv\; \neg \mathrm{AF}\,\neg\varphi
\;\equiv\; \ER{\bot}{\varphi}, $ and $
\mathrm{AG}\,\varphi \;\equiv\; \neg \mathrm{EF}\,\neg\varphi
\;\equiv\; \AR{\bot}{\varphi}.$
Their time-bounded counterparts are defined analogously: $\mathrm{EF}_{\kI}\,\varphi \;\equiv\; \tEU{\top}{\varphi}{\kI},
\mathrm{AF}_{\kI}\,\varphi \;\equiv\; \tAU{\top}{\varphi}{\kI}, 
\mathrm{EG}_{\kI}\,\varphi \;\equiv\;
\tER{\bot}{\varphi}{\kI},$ and $
\mathrm{AG}_{\kI}\,\varphi \;\equiv\;
\tAR{\bot}{\varphi}{\kI}.
$

%Hence, the core temporal operators of the logic are $\mathbf{U}$, $\mathbf{R}$, their time-bounded variants, and the next-step modalities. All standard CTL and TCTL operators can be derived from this basis.


%In order to define the syntax and semantics of TCTL, we first introduce the underlying models and the syntax of formulas. We consider a standard clocked semantics for TCTL, where the models are timed transition systems (TTS) with explicit clock valuations. A {clock} is a positive real-valued variable that increases uniformly over time. The semantics of TCTL formulas is defined in terms of paths in the combined discrete--time semantics of TTS, where delay steps advance all clocks uniformly. 

%We interpret TCTL formulas over timed transition systems (TTS) defined as follows.
The semantics of TCTL formulas is defined over timed transition systems (TTS), which are a standard model for real-time systems.

\begin{definition}[Timed Transition System (TTS)]\label{def:tts}
A {Timed Transition System} over a finite clock set $\clockSet$ is a tuple $\mathcal{T}
=
\langle S, s_0, \mathsf{Inv}, E, \mathcal{P}, L \rangle
$ where $S$ is a (possibly infinite) set of {locations}, $s_0 \in S$ the initial location, $\mathsf{Inv}$ assigns to each location a {location invariant} (a zone over $\clockSet$), $E \subseteq S \times 2^{\clockSet} \times S$ a set of {discrete edges} each carrying the set of clocks it resets, $\mathcal{P}$ a finite set of atomic propositions, and $L : S \to 2^{\mathcal{P}}$ a labeling function.
\end{definition}

A {configuration} of $\mathcal{T}$ is a pair $(s,\nu)$ with $s \in S$, $\nu \in \clockValuations{\clockSet}$ and $\nu\models\mathsf{Inv}(s)$; we write $\mathrm{Conf}(\mathcal{T})$ for the set of configurations. The satisfaction relation ranges over configurations, and the operational behaviour is given by two relations on configurations:
\begin{itemize}
  \item {delay} steps $(s,\nu) \xrightarrow{\;t\;}_{\tau} (s,\nu+t)$ for $t \in \mathbb{R}_{\ge 0}$ such that $\nu+t'\models\mathsf{Inv}(s)$ for all $t'\in[0,t]$ (the invariant holds throughout the delay; by convexity of zones it suffices that $\nu$ and $\nu+t$ satisfy it), where $(\nu+t)(x) = \nu(x)+t$; and
  \item {discrete} steps $(s,\nu) \rightarrow_d (s',\nu[R:=0])$ for every edge $(s,R,s') \in E$ with $\nu\models\mathsf{Inv}(s)$ and $\nu[R:=0]\models\mathsf{Inv}(s')$, where $\nu[R:=0](x) = 0$ if $x \in R$ and $\nu(x)$ otherwise.
\end{itemize}
Location invariants are what let a TTS express {urgency}: by bounding how long time may elapse in a location they can force a discrete step before a deadline, which is indispensable for satisfying universal time-bounded eventualities such as $\mathrm{AF}_{[1,2]}\,p$ (an invariant-free model would admit a run lingering at one location across the whole window). We assume $\mathcal{T}$ is {non-blocking}: from every configuration there {exists} a non-Zeno run (this existence is the load-bearing requirement; informally it holds when time may diverge within the current invariant, or a discrete step remains available as the invariant tightens). In particular $\mathcal{T}$ has no timelocks, so a run with finitely many discrete steps may always be completed by an unbounded final delay. Delay steps satisfy time additivity: $(s,\nu) \xrightarrow{t}_{\tau} (s,\nu+t)$ and $(s,\nu+t)\xrightarrow{u}_{\tau}(s,\nu+t+u)$ imply $(s,\nu)\xrightarrow{t+u}_{\tau}(s,\nu+t+u)$.

A {run} from a configuration $c_0=(s_0,\nu_0)$ is a maximal sequence of delay and discrete steps
\[
\rho \;=\; (s_0,\nu_0)\xrightarrow{t_0}_{\tau} \rightarrow_d (s_1,\nu_1)\xrightarrow{t_1}_{\tau} \rightarrow_d (s_2,\nu_2)\cdots,\qquad t_i \in \mathbb{R}_{\ge 0},
\]
with finitely or infinitely many discrete steps; a run with finitely many discrete steps ends in an unbounded final delay $(s_k,\nu_k)\xrightarrow{\infty}_{\tau}$ (pure time elapse). A {position} of $\rho$ is a pair $\pi=(i,\delta)$ with $0\le\delta\le t_i$ (and $0\le\delta<\infty$ on a final unbounded delay), denoting the configuration $\rho(\pi)=(s_i,\nu_i+\delta)$ and carrying {elapsed time} $\mathit{time}(\pi)=T_i+\delta$ where $T_i=\sum_{j<i} t_j$. Positions are ordered lexicographically: $(i,\delta)\preceq(i',\delta')$ iff $i<i'$, or $i=i'$ and $\delta\le\delta'$. Quantifying over positions, rather than over bare absolute time, keeps the configuration unambiguous at discrete jumps, where the pre-jump position $(i,t_i)$ and the post-jump position $(i{+}1,0)$ share the same elapsed time but denote distinct configurations.

All steps of a run respect the invariants (that is what makes a run a sequence of configurations). The run is {Zeno} if it has infinitely many discrete steps yet $\sum_i t_i<\infty$, and {time-divergent} ({non-Zeno}) if $\sum_i t_i=\infty$; in particular every run ending in an unbounded delay is non-Zeno. We write $\mathrm{Runs}^{\infty}(c)$ for the set of non-Zeno runs from $c$; by non-blockingness it is non-empty for every configuration $c$, so the modalities below are never vacuous and $\mathrm{EF}$ is reflexive. Evaluating the temporal modalities over non-Zeno runs only is the standard dense-time semantics of~\cite{AlurCourcoubetisDill1990}, whose path quantifiers already range over time-divergent paths; the loop-acceptance condition of Section~\ref{sec:tableau} is what makes the tableau honour it, and Section~\ref{sec:proofs} shows the restriction can be lifted by dropping that condition.

The satisfaction relation $\mathcal{T},(s,\nu)\models\varphi$, read ``$\varphi$ holds at configuration $(s,\nu)$,'' is defined by induction on $\varphi$.

\begin{definition}[Satisfaction]\label{def:sat}
The Boolean cases are standard: $(s,\nu)\models\top$ always; $(s,\nu)\not\models\bot$; $(s,\nu)\models p$ iff $p\in L(s)$; $(s,\nu)\models\neg p$ iff $p\notin L(s)$; and $\land,\lor$ as usual. For the timed modalities, with $\kI$ a time bound and the untimed operators recovered as the case $\kI=[0,\infty)$:
\begin{align*}
(s,\nu)\models \tEU{\varphi}{\psi}{\kI}
&\iff \exists\,\rho\in\mathrm{Runs}^\infty(s,\nu)\ \exists\,\text{a position }\pi\text{ of }\rho:\ \mathit{time}(\pi)\in \kI,\ \rho(\pi)\models\psi,\ \text{and}\ \forall\,\pi'\prec\pi:\ \rho(\pi')\models\varphi,\\
(s,\nu)\models \tAU{\varphi}{\psi}{\kI}
&\iff \forall\,\rho\in\mathrm{Runs}^\infty(s,\nu)\ \exists\,\text{a position }\pi\text{ of }\rho:\ \mathit{time}(\pi)\in \kI,\ \rho(\pi)\models\psi,\ \text{and}\ \forall\,\pi'\prec\pi:\ \rho(\pi')\models\varphi.
\end{align*}
Quantifying over positions $\pi$ (with the lexicographic order $\prec$) rather than over absolute time makes the clauses unambiguous at discrete jumps. Release is the De~Morgan dual of until, which keeps NNF closed under negation:
\[
\tER{\varphi}{\psi}{\kI}\;\equiv\; \neg\,\tAU{\neg\varphi}{\neg\psi}{\kI},\qquad
\tAR{\varphi}{\psi}{\kI}\;\equiv\; \neg\,\tEU{\neg\varphi}{\neg\psi}{\kI}.
\]
Unfolding the dual gives the positive reading used by the tableau: $(s,\nu)\models \tAR{\varphi}{\psi}{\kI}$ iff for {every} $\rho\in\mathrm{Runs}^\infty(s,\nu)$ and every position $\pi$ of $\rho$ with $\mathit{time}(\pi)\in \kI$, either $\rho(\pi)\models\psi$, or $\rho(\pi')\models\varphi$ for some position $\pi'\prec\pi$; and $\tER{\varphi}{\psi}{\kI}$ is the same with $\exists\,\rho$ in place of $\forall\,\rho$.
\end{definition}


\subsection{Zones}
As introduced above, the semantics of TCTL is defined over configurations $(s,\nu)$, where $s$ is a location and $\nu$ is a clock valuation.
%The semantics introduced above is defined over concrete configurations $(s,\nu)$.
In the tableau construction, however, we reason symbolically about
sets of clock valuations rather than individual valuations.
To obtain a finite and effectively manipulable representation
of such sets, we use zones, i.e., convex sets of valuations
definable by conjunctions of difference constraints \cite{Bouyer05TimedAutomata,AlurDill1994}.
Zones form a class closed under the operations required
by the clocked semantics of TCTL, namely intersection,
time elapse, and clock reset.

\begin{definition}[Zones]\label{def:zones}
A {zone} is a set $Z \subseteq \mathbb{R}_{\ge 0}^{\clockSet}$
definable by a finite conjunction of {difference constraints}
of the form $x - y \prec c$,
where $x,y \in \clockSet \cup \{0\}$ and $(c,\prec)$ ranges over $\big(\mathbb{Z}\times\{<,\le\}\big)\cup\{(\infty,<)\}$ (the value $\infty$ pairs only with $<$, denoting the absence of a bound).
\end{definition}

Crucially, each constraint carries a {strictness flag} $\prec\,\in\{<,\le\}$, so that zones can represent both closed and open bounds; this is mandatory because the (non-punctual) time bounds of Definition~\ref{def:time_bounds} include open intervals such as $(a,b)$, and the rules of Section~\ref{sec:tableau} intersect zones with such bounds. Concretely, zones are stored as {difference-bound matrices} (DBMs) over the bound domain $\big(\mathbb{Z}\times\{<,\le\}\big)\cup\{(\infty,<)\}$, where $(\infty,<)$ is the top element $\langle$no constraint$\rangle$~\cite{Bouyer05TimedAutomata,AlurDill1994}. In the definition above we distinguish the symbol $0$, a fictitious clock constantly equal to $0$; thus a clock bound $x \prec c$ (resp.\ $x \succ c$) is written $x - 0 \prec c$ (resp.\ $0 - x \prec -c$).

A zone is {empty} if it contains no valuation, denoted $\emptyset$; emptiness is detected as a negative cycle in the DBM. The {universal} zone is $\Zuniv = \mathbb{R}_{\ge 0}^{\clockSet}$. The operations the tableau requires are the following, and each maps zones to zones (so the class is effectively closed under them):
\begin{itemize}
  \item {Intersection} $Z_1\cap Z_2$ is set-theoretic intersection (conjunction of the two constraint systems).
  \item {Guard intersection} $Z\cap(x\in \kI)$ with a time bound $\kI$ conjoins the (possibly strict) bounds of $\kI$ on the single clock $x$; this is the operation that requires the strictness flag.
  \item {Reset} of $R\subseteq\clockSet$ is $Z[R := 0] = \{\, \nu' \mid \exists \nu \in Z :\nu'(x)=0 \text{ if } x\in R,\text{ and } \nu'(x)=\nu(x) \text{ otherwise} \,\}$, computed by removing and re-zeroing the rows/columns of $R$.
  \item {Time-successor} (upward closure) $\uparrow Z = \{\nu + t \mid \nu \in Z,\ t \in \mathbb{R}_{\ge 0}\}$, computed by dropping every upper bound $x-0\prec c$ while keeping all differences $x-y\prec c$ and lower bounds $0-x\prec c$ (the fictitious clock $0$ does not advance).
\end{itemize}
We remark that resetting a clock by the {constraint} $Z\cap(x=0)$ coincides with the {reset} operation $Z[\{x\}:=0]$ exactly when $x$ is unconstrained in $Z$ (which is the case at clock allocation, Section~\ref{sec:tableau}); the equality $x=0$ used there is a zone constraint and is unrelated to the ban on punctual {modal} bounds of Definition~\ref{def:time_bounds}.

To keep the set of reachable zones finite we compare them modulo {max-constant extrapolation}. Let $M\in\mathbb{N}$ be the largest constant occurring in any time bound of the input formula.
\begin{definition}[$M$-equivalence of zones]\label{def:zone_equiv}
Let $\alpha_M$ be the standard max-constant extrapolation, which relaxes (weakens) to $\infty$ every upper-bound constant exceeding $M$ and weakens to $>M$ every lower bound above $M$, so that all differences beyond $M$ are identified. For zones $Z_1,Z_2$, write $Z_1 \cong_M Z_2 \iff \alpha_M(Z_1)=\alpha_M(Z_2)$, comparing the {canonical} (shortest-path-closed) DBMs of the extrapolated zones.
\end{definition}

We assume each zone is stored as $\mathrm{canonical}(\alpha_M(Z))$, the canonical DBM of its extrapolation; this is a unique representative of its $\cong_M$-class, so $\cong_M$-equivalent zones have identical representations and are compared by syntactic equality. Soundness of this abstraction relies on a property of our rules: every clock guard {tested} by the tableau (Section~\ref{sec:tableau}) is a single-clock bound $x\in\kI$, i.e.\ the construction is {diagonal-free} in its guards. Max-constant extrapolation is sound for diagonal-free guards even though reachable zones may contain diagonal constraints $x-y\prec c$~\cite{Bouyer05TimedAutomata}; the diagonal entries are inert for branch closure and only record the relative offsets of concurrently active clocks.



